\documentclass{article}
\usepackage{amsmath,amssymb}
\usepackage{xurl}
\usepackage[hidelinks]{hyperref}
% Shared commands for the notes in docs/paper-gaps/.
% Every note loads this file via % Shared commands for the notes in docs/paper-gaps/.
% Every note loads this file via % Shared commands for the notes in docs/paper-gaps/.
% Every note loads this file via \input{command}. Project-wide notation
% belongs here, never in per-note redefinitions.

% --- Project configuration ------------------------------------------------
\newcommand{\projectIssueBase}{https://github.com/LionSR/QICLean/issues/}
\newcommand{\projectPrBase}{https://github.com/LionSR/QICLean/pull/}
\newcommand{\projectDocBase}{https://sirui-lu.com/QICLean/docs/}
% ---------------------------------------------------------------------------

\DeclareUnicodeCharacter{2080}{\ifmmode{}_{0}\else\textsubscript{0}\fi}
\DeclareUnicodeCharacter{2081}{\ifmmode{}_{1}\else\textsubscript{1}\fi}
\DeclareUnicodeCharacter{2082}{\ifmmode{}_{2}\else\textsubscript{2}\fi}
\DeclareUnicodeCharacter{2099}{\ifmmode{}_{n}\else\textsubscript{n}\fi}

\newcommand{\Lean}{\textsc{Lean}}

% A formal declaration name, upright sans, breakable at underscores.
\ExplSyntaxOn
\NewDocumentCommand{\leanid}{m}{
  \ifmmode
    \textsf{\detokenize{#1}}
  \else
    \group_begin:
    \urlstyle{sf}
    \tl_set:Nn \l_tmpa_tl {#1}
    \tl_replace_all:Nnn \l_tmpa_tl {\_} {_}
    \exp_args:NV \path \l_tmpa_tl
    \group_end:
  \fi
}
\ExplSyntaxOff

% Traceability links; use in footnotes, not in the running prose.
\newcommand{\ghissue}[1]{\href{\projectIssueBase#1}{GitHub issue \##1}}
\newcommand{\ghpr}[1]{\href{\projectPrBase#1}{PR \##1}}
\newcommand{\doclink}[1]{\href{\projectDocBase#1.html}{\path{#1}}}

% Domain notation, using \providecommand so a note-local refinement stays
% possible.
\providecommand{\tr}{\operatorname{tr}}

% Dirac notation and the identity operator, as in blueprint/src/macros.
\usepackage{dsfont}
\providecommand{\ket}[1]{|#1\rangle}
\providecommand{\bra}[1]{\langle #1|}
\providecommand{\braket}[2]{\langle #1 | #2 \rangle}
\providecommand{\ketbra}[2]{|#1\rangle\!\langle #2|}
\providecommand{\Id}{\mathds{1}}
\providecommand{\one}{\mathds{1}}
\providecommand{\C}{\mathbb{C}}
\providecommand{\MN}[1]{M_{#1}(\C)}
\providecommand{\MD}{M_D(\C)}

% Verdict marker: \gapnote{<kind>}{<status>}. Kind is the policy
% classification (clarification, local-correction, scope-restriction,
% unfaithful, false-source, open-gap); status is open, wip (elimination
% underway), resolved, or historical. Machine-read by the site index;
% prints a verdict line.
\newcommand{\gapnote}[2]{\par\noindent\textbf{Verdict:} #1 (#2).\par}
. Project-wide notation
% belongs here, never in per-note redefinitions.

% --- Project configuration ------------------------------------------------
\newcommand{\projectIssueBase}{https://github.com/LionSR/QICLean/issues/}
\newcommand{\projectPrBase}{https://github.com/LionSR/QICLean/pull/}
\newcommand{\projectDocBase}{https://sirui-lu.com/QICLean/docs/}
% ---------------------------------------------------------------------------

\DeclareUnicodeCharacter{2080}{\ifmmode{}_{0}\else\textsubscript{0}\fi}
\DeclareUnicodeCharacter{2081}{\ifmmode{}_{1}\else\textsubscript{1}\fi}
\DeclareUnicodeCharacter{2082}{\ifmmode{}_{2}\else\textsubscript{2}\fi}
\DeclareUnicodeCharacter{2099}{\ifmmode{}_{n}\else\textsubscript{n}\fi}

\newcommand{\Lean}{\textsc{Lean}}

% A formal declaration name, upright sans, breakable at underscores.
\ExplSyntaxOn
\NewDocumentCommand{\leanid}{m}{
  \ifmmode
    \textsf{\detokenize{#1}}
  \else
    \group_begin:
    \urlstyle{sf}
    \tl_set:Nn \l_tmpa_tl {#1}
    \tl_replace_all:Nnn \l_tmpa_tl {\_} {_}
    \exp_args:NV \path \l_tmpa_tl
    \group_end:
  \fi
}
\ExplSyntaxOff

% Traceability links; use in footnotes, not in the running prose.
\newcommand{\ghissue}[1]{\href{\projectIssueBase#1}{GitHub issue \##1}}
\newcommand{\ghpr}[1]{\href{\projectPrBase#1}{PR \##1}}
\newcommand{\doclink}[1]{\href{\projectDocBase#1.html}{\path{#1}}}

% Domain notation, using \providecommand so a note-local refinement stays
% possible.
\providecommand{\tr}{\operatorname{tr}}

% Dirac notation and the identity operator, as in blueprint/src/macros.
\usepackage{dsfont}
\providecommand{\ket}[1]{|#1\rangle}
\providecommand{\bra}[1]{\langle #1|}
\providecommand{\braket}[2]{\langle #1 | #2 \rangle}
\providecommand{\ketbra}[2]{|#1\rangle\!\langle #2|}
\providecommand{\Id}{\mathds{1}}
\providecommand{\one}{\mathds{1}}
\providecommand{\C}{\mathbb{C}}
\providecommand{\MN}[1]{M_{#1}(\C)}
\providecommand{\MD}{M_D(\C)}

% Verdict marker: \gapnote{<kind>}{<status>}. Kind is the policy
% classification (clarification, local-correction, scope-restriction,
% unfaithful, false-source, open-gap); status is open, wip (elimination
% underway), resolved, or historical. Machine-read by the site index;
% prints a verdict line.
\newcommand{\gapnote}[2]{\par\noindent\textbf{Verdict:} #1 (#2).\par}
. Project-wide notation
% belongs here, never in per-note redefinitions.

% --- Project configuration ------------------------------------------------
\newcommand{\projectIssueBase}{https://github.com/LionSR/QICLean/issues/}
\newcommand{\projectPrBase}{https://github.com/LionSR/QICLean/pull/}
\newcommand{\projectDocBase}{https://sirui-lu.com/QICLean/docs/}
% ---------------------------------------------------------------------------

\DeclareUnicodeCharacter{2080}{\ifmmode{}_{0}\else\textsubscript{0}\fi}
\DeclareUnicodeCharacter{2081}{\ifmmode{}_{1}\else\textsubscript{1}\fi}
\DeclareUnicodeCharacter{2082}{\ifmmode{}_{2}\else\textsubscript{2}\fi}
\DeclareUnicodeCharacter{2099}{\ifmmode{}_{n}\else\textsubscript{n}\fi}

\newcommand{\Lean}{\textsc{Lean}}

% A formal declaration name, upright sans, breakable at underscores.
\ExplSyntaxOn
\NewDocumentCommand{\leanid}{m}{
  \ifmmode
    \textsf{\detokenize{#1}}
  \else
    \group_begin:
    \urlstyle{sf}
    \tl_set:Nn \l_tmpa_tl {#1}
    \tl_replace_all:Nnn \l_tmpa_tl {\_} {_}
    \exp_args:NV \path \l_tmpa_tl
    \group_end:
  \fi
}
\ExplSyntaxOff

% Traceability links; use in footnotes, not in the running prose.
\newcommand{\ghissue}[1]{\href{\projectIssueBase#1}{GitHub issue \##1}}
\newcommand{\ghpr}[1]{\href{\projectPrBase#1}{PR \##1}}
\newcommand{\doclink}[1]{\href{\projectDocBase#1.html}{\path{#1}}}

% Domain notation, using \providecommand so a note-local refinement stays
% possible.
\providecommand{\tr}{\operatorname{tr}}

% Dirac notation and the identity operator, as in blueprint/src/macros.
\usepackage{dsfont}
\providecommand{\ket}[1]{|#1\rangle}
\providecommand{\bra}[1]{\langle #1|}
\providecommand{\braket}[2]{\langle #1 | #2 \rangle}
\providecommand{\ketbra}[2]{|#1\rangle\!\langle #2|}
\providecommand{\Id}{\mathds{1}}
\providecommand{\one}{\mathds{1}}
\providecommand{\C}{\mathbb{C}}
\providecommand{\MN}[1]{M_{#1}(\C)}
\providecommand{\MD}{M_D(\C)}

% Verdict marker: \gapnote{<kind>}{<status>}. Kind is the policy
% classification (clarification, local-correction, scope-restriction,
% unfaithful, false-source, open-gap); status is open, wip (elimination
% underway), resolved, or historical. Machine-read by the site index;
% prints a verdict line.
\newcommand{\gapnote}[2]{\par\noindent\textbf{Verdict:} #1 (#2).\par}

\title{Asymptotic Conversion by Protocols without Classical Feedforward}
\author{}
\date{2026-11-18}
\begin{document}
\maketitle
\gapnote{scope-restriction}{open}

\section{The source relation}
In the paragraph ``Phases of matter'' of~\cite{Piroli2021LOCC}, a sequence
$\{\ket{\psi_N}\}$ is convertible into $\{\ket{\varphi_N}\}$ if there is a
fixed number $k$ of compositions of channels assisted by measurements and
classical communication, each of depth $f(N)$, whose outputs $\sigma_N$
satisfy
\begin{align}
    \lVert\sigma_N-\ket{\varphi_N}\bra{\varphi_N}\rVert_1 & \longrightarrow 0.
    \label{eq:psc21_local_channel_phase_scope_source_convergence}
\end{align}
Here $f$ is polylogarithmic. Mutual convertibility defines the source's
phase equivalence. The fixed number of compositions and the polylogarithmic
depth bound are allowed to depend on the pair of sequences, but not on $N$.

\section{The enlarged-site protocol model}
The protocols are unbranched circuits on a geometry, whose layers carry
trace-preserving channels on pairwise disjoint cells and trace-preserving
channels at the remaining sites. They act on arbitrary input operators and
allow changes of local dimension. A layer using a cell of positive cost costs
one depth unit, regardless of the dimensions of the bases it acts on.
Intermediate dimensions may depend on the system size without a uniform
bound. The model has no classical feedforward between separated sites. For
this class, approximate conversion means
$\lVert\Psi(\rho)-\sigma\rVert_1\le\varepsilon$, with protocol depth at most
$T$.\footnote{Formal declaration:
    \leanid{QuantumCircuit.IsApproxConversion} in
    \doclink{QICLean/Circuit/Model/Approximate}.}
Composition adds the depths and errors:
\begin{align}
    \lVert\Psi_2(\Psi_1(\rho))-\tau\rVert_1
     & \le \lVert\Psi_2(\Psi_1(\rho)-\sigma)\rVert_1
    +\lVert\Psi_2(\sigma)-\tau\rVert_1               \\
     & \le \lVert\Psi_1(\rho)-\sigma\rVert_1
    +\lVert\Psi_2(\sigma)-\tau\rVert_1
    \le\varepsilon_1+\varepsilon_2.
    \label{eq:psc21_local_channel_phase_scope_composition_error}
\end{align}
The middle inequality is the contraction of the trace norm under a
trace-preserving completely positive map, which holds for every input
operator, Hermitian or not.\footnote{Formal declarations:
    \leanid{KrausFamily.IsTP.rectangularTraceNorm_map_le} in
    \doclink{QICLean/Channel/KrausFamilyTraceNorm} and
    \leanid{QuantumCircuit.IsApproxConversion.trans}.}

On the ring of $N$ sites, whose cells are the pairs $\{k,k+1\}$ of cost $1$, the
$T$-fold neighbourhood of a region is its ball of radius $T$. A state converted in
depth $T$ from a product of unit-trace factors therefore has
$\operatorname{tr}(\sigma AB)=\operatorname{tr}(\sigma A)\operatorname{tr}(\sigma B)$
for $A$ and $B$ supported on regions at ring distance larger than $2T$. This is the
source's necessary condition, Proposition~\texttt{propQCA2}, in the enlarged-site
model; the source states it for pure states prepared by unitary circuits.\footnote{Formal
    declarations: \leanid{QuantumCircuit.Geometry.ring} in
    \doclink{QICLean/Circuit/Geometry/Ring} and
    \leanid{QuantumCircuit.IsConversion.trace_mul_mul_eq_of_isRingSeparated} in
    \doclink{QICLean/Circuit/Model/Conversion}.}

For families on a family of geometries indexed by a size $N$, the error tends
to zero and the depth is eventually bounded by $C(1+\log(N+1))^p$, for a
nonnegative constant $C$ and an integer $p\ge0$. Mutual conversion is an
equivalence relation on families of operators.\footnote{Formal declarations:
    \leanid{QuantumCircuit.IsAsymptoticConversion} and
    \leanid{QuantumCircuit.isPhaseEquivalent_equivalence} in
    \doclink{QICLean/Circuit/Model/Approximate}.}
The shift in the logarithm only affects the bound at small sizes. Finite
compositions are already included: their depths add, and a finite sum of
polylogarithmic bounds is polylogarithmic. No bound on the dimension of
onsite ancillas is imposed.

This gate convention differs from the source. In the Supplement's proof of
Proposition~2 (p.~8), the intersite gates act only on the physical qudits;
the entanglement bound across a cut depends on their fixed dimension $d$,
not on the number of onsite ancillas. Allowing one intersite layer to act
on an arbitrarily enlarged local space does not by itself give a
uniform-depth simulation in that convention. No such simulation or
inclusion between the two asymptotic conversion relations is proved here.

\section{Remaining scope}
The equivalence proved here belongs to the enlarged-site protocol model. It
is not identified with the source's phase equivalence or asserted to be a
subclass at the same asymptotic depth. The classification asserting that
every translation-invariant matrix product state belongs to the trivial
source phase is not obtained. Circuits whose later layers depend on the
outcomes of earlier ones are available as adaptive protocols, with additive
depth under composition, but they are not included in the asymptotic
relation formalized here.

The remaining source comparison needs a decomposition into a fixed number
of \textit{QCcc} blocks, with that number independent of system size. Each
block must satisfy the source's internal measurement and control order and
its rule of one measurement per site. The comparison must also distinguish
fixed-dimensional physical qudits from onsite ancillas and control the
intersite gate cost uniformly in the system size. A proved simulation bound
or explicit restrictions on intermediate dimensions could supply this
resource comparison; the encoding of a protocol with bounded bases into a
fixed alphabet at every site, of the same depth, is one such ingredient. Every
protocol of the enlarged-site model is a preparation of basis vectors on
environments at the sites, an unbranched circuit of unitary gates of no larger
depth on the same geometry, and a partial trace over the output environments.
The environments lie in the local bases, and their dimensions depend on the
protocol, so this does not yet act on fixed-dimensional physical qudits, which
requires encoding the environments into qudits and routing them.\footnote{Formal declaration:
    \leanid{QuantumCircuit.IsProtocol.isUnitaryDilation} in
    \doclink{QICLean/Circuit/Compile/Dilation}.} Once
those source-model depth bounds are established, the contraction of the
trace norm supplies the same error estimate.

\bibliographystyle{alpha}
\bibliography{references}
\end{document}
